% =====================================================================
% EJE B - LAS TECNICAS DEL TRACK A  ---  RESPONSABLE: INTEGRANTE 2
% Aqui se justifica el track CON LITERATURA, no con opinion.
% De cada familia: fortalezas, limites y SUPUESTOS bajo los que funciona.
% =====================================================================
\subsection{Axis B: Techniques}
\label{sec:ejeB}

% TODO Int.2 - subtemas sugeridos:
%  (a) Gradient boosting vs redes neuronales en datos tabulares
%  (b) El problema es RANKING dentro de un grupo, no clasificacion
%      binaria: learning to rank, conditional logit, Plackett-Luce
%  (c) Calibracion de probabilidades: Platt, isotonica
%  (d) Explicabilidad: SHAP / permutation importance

% --- LINEA DE PRUEBA: borrar cuando Int.2 escriba la seccion real ---
